\section{Validation}
\label{sec:validation}

We validate the four theorems on simulated phenotyping cohorts with
known ground-truth clinical state. All results are reproducible from
\texttt{scripts/run.R} (seed \texttt{20260521}).

\subsection{Data-generating process}

Each cohort has $n$ patients. The true clinical state $\Y_i \in
\{0,1\}$ is Bernoulli with prevalence $\pi$. Latent severity is $S_i
\sim \mathrm{Normal}(\mu_{S,\Y_i}, 1)$ with $\mu_{S,0} = 0$ and
$\mu_{S,1} = 1.5$, so true cases are more severe on average. The
diagnosis code is generated by the coding mechanism
\eqref{eq:trans-coding}: for a true case the code probability is
$\mathrm{logit}^{-1}(b_0 + b_1 S_i)$, and for a non-case it is the
false-positive rate $\mathrm{fpr}$. Unless stated otherwise $b_0 = 0$,
$\mathrm{fpr} = 0.05$ (code specificity $0.95$), and the coding slope
is $b_1 = 1$, the informative-coding regime in which code sensitivity
rises with severity. Three estimators of $\pi$ are compared: code-only
($\hat\pi = \bar C$), chart-review-calibrated (adjudicate a subsample
of $m$ patients, fit the coding model on the subsample, deconvolve
cohort prevalence via \eqref{eq:deconvolve}), and oracle (full chart
review, $\hat\pi = \bar\Y$).

\subsection{Glass ceiling (\cref{thm:glass-ceiling})}

We generate one cohort with $n = 50{,}000$, $\pi = 0.12$, and $b_1 =
1$. The true marginal code sensitivity is $0.776$ and specificity is
$0.950$. The observed code frequency is $\bar C = 0.1377$: the code is
positive in $13.8\%$ of patients while only $12\%$ have the condition.

We then enumerate the solution surface of \cref{thm:glass-ceiling}: the
set of $(\pi, \mathrm{sens}, \mathrm{spec})$ triples that reproduce the
observed code frequency exactly. On a grid of prevalence and
specificity values, $160$ admissible triples reproduce $\bar C =
0.1377$, with implied prevalence ranging from $0.06$ to $0.40$
(\Cref{fig:glass}A). Every triple matches the observed code frequency
to within $3 \times 10^{-17}$, machine precision. The true prevalence
$0.12$ is one point on this surface, and no feature of the code data
distinguishes it from a model with prevalence $0.40$. This is the glass
ceiling: code data alone supplies one scalar constraint on three
unknowns.

\begin{figure}[t]
\centering
\includegraphics[width=0.95\linewidth]{glass_ceiling.pdf}
\caption{(A) Glass-ceiling solution surface. Each curve fixes a code
specificity; every point reproduces the observed code frequency $\bar
C = 0.138$ exactly. The true parameters (red) are one point among a
continuum, with implied prevalence spanning $0.06$ to $0.40$. (B)
Chart review restores identifiability: RMSE of the
chart-review-calibrated prevalence estimate falls as the reviewed
subsample size $m$ grows, while the code-only estimator's error is
flat (it is bias, not sampling error).}
\label{fig:glass}
\end{figure}

\subsection{Identifiability with chart review (\cref{thm:identifiability-chart})}

We test whether a representative chart-reviewed subsample recovers
prevalence. For $\pi = 0.12$, $b_1 = 1$, and $n = 20{,}000$, we draw a
simple random subsample of size $m$, adjudicate the true state, fit the
coding model on the subsample, and deconvolve cohort prevalence. Over
$200$ replicates per subsample size, the calibrated estimator's RMSE
decreases monotonically with $m$: $0.097$ at $m = 50$, $0.039$ at $m =
100$, $0.025$ at $m = 200$, $0.014$ at $m = 500$, $0.0102$ at $m =
1{,}000$, and $0.0073$ at $m = 2{,}000$ (\Cref{tab:chart},
\Cref{fig:glass}B). The calibrated estimator's bias is small at every
$m$ (magnitude below $0.013$ at $m = 50$ and below $0.001$ from $m =
1{,}000$ on); the RMSE at small $m$ is dominated by the variance of
fitting the coding model on few adjudicated cases.

The code-only estimator, by contrast, has RMSE essentially constant
near $0.0177$ across all $m$, because its error is bias, not sampling
noise: it estimates the code frequency, not the prevalence, and chart
review of patients it never uses cannot help it.
\Cref{thm:identifiability-chart} is confirmed: a chart-reviewed
subsample covering both classes identifies the coding model, and cohort
prevalence becomes identifiable, with estimation error shrinking at the
usual $m^{-1/2}$ rate.

\begin{table}[ht]
\centering
\begin{tabular}{rrrrr}
\toprule
$m$ & calibrated RMSE & calibrated bias & code-only RMSE & code-only bias \\
\midrule
$50$    & $0.0971$ & $+0.0129$  & $0.0178$ & $+0.0177$ \\
$100$   & $0.0389$ & $-0.0002$  & $0.0178$ & $+0.0176$ \\
$200$   & $0.0255$ & $+0.0036$  & $0.0179$ & $+0.0177$ \\
$500$   & $0.0144$ & $+0.0013$  & $0.0172$ & $+0.0171$ \\
$1{,}000$ & $0.0102$ & $+0.0010$ & $0.0178$ & $+0.0177$ \\
$2{,}000$ & $0.0073$ & $-0.0000$ & $0.0177$ & $+0.0175$ \\
\bottomrule
\end{tabular}
\caption{Chart-review-calibrated versus code-only prevalence
estimation, $\pi = 0.12$, $200$ replicates per row. The calibrated
RMSE shrinks with the reviewed subsample size $m$; the code-only error
is a fixed bias.}
\label{tab:chart}
\end{table}

\subsection{Code-frequency consistency (\cref{thm:code-total})}

The theorem predicts that a fitted phenotype model reproduces the
empirical code frequency exactly. Over $200$ replicates with $n =
20{,}000$ and a chart-review subsample of $m = 1{,}000$, the
calibrated triple $(\hat\pi, \widehat{\mathrm{sens}},
\widehat{\mathrm{spec}})$ has code-frequency residual $|\,\hat q -
\bar C\,|$ with median $0$ and maximum $3 \times 10^{-17}$, machine
precision.

The sharper check is the one \cref{thm:code-total} warns about. We take
a model with a \emph{deliberately wrong} prevalence (twice the truth)
and re-solve its sensitivity to match the observed code frequency. This
wrong-prevalence model reproduces $\bar C$ exactly: median and maximum
residual both $0$. A marginal-fit check, comparing predicted to
observed coding rate, cannot tell the wrong-prevalence model from the
correct one. The diagnostic practitioners reach for is uninformative
about the quantity of interest, exactly as the theorem states.

\subsection{Bias under informative coding (\cref{thm:bias-informative})}

\paragraph{The coding slope sweep.} We sweep the coding-mechanism
severity slope $b_1 \in \{0, 0.5, 1, 1.5, 2, 3\}$, $\pi = 0.12$, $200$
replicates per value. At $b_1 = 0$ the coding mechanism does not depend
on severity, C2 holds, and the severity-coding correlation $\rho$ is
$0$; as $b_1$ grows the marginal code sensitivity rises (from $0.50$ at
$b_1 = 0$ to $0.90$ at $b_1 = 3$) because severe cases are
preferentially coded (\Cref{tab:informative}).

The code-only estimator's bias tracks the C2 violation. At $b_1 = 0$
the bias is $-0.016$: with sensitivity only $0.50$, half of true cases
are missed, and the miss deficit $\pi(1 - \mathrm{sens})$ slightly
outweighs the false-positive inflation $(1-\pi)\,\mathrm{fpr}$, so
code-only prevalence \emph{understates} the truth. As $b_1$ rises the
sensitivity deficit shrinks faster than the false-positive inflation,
and the code-only bias swings positive: $+0.004$ at $b_1 = 0.5$,
$+0.017$ at $b_1 = 1$, $+0.025$ at $b_1 = 1.5$, $+0.028$ at $b_1 = 2$,
and $+0.032$ at $b_1 = 3$. The magnitude of the bias is governed by the
coding mechanism, not by the cohort size, exactly as
\cref{thm:bias-informative} predicts. The chart-review-calibrated
estimator stays near-unbiased throughout: its bias magnitude is below
$0.002$ at every $b_1$, and its RMSE falls from $0.019$ at $b_1 = 0$ to
$0.0084$ at $b_1 = 3$ (\Cref{fig:informative}A).

\begin{table}[ht]
\centering
\begin{tabular}{rrrrr}
\toprule
$b_1$ & $\rho$ & code-only bias & calibrated bias & MCSE \\
\midrule
$0.0$ & $0.000$ & $-0.0161$ & $+0.0030$ & $0.0013$ \\
$0.5$ & $0.994$ & $+0.0045$ & $+0.0008$ & $0.0009$ \\
$1.0$ & $0.954$ & $+0.0174$ & $+0.0010$ & $0.0008$ \\
$1.5$ & $0.894$ & $+0.0247$ & $+0.0015$ & $0.0007$ \\
$2.0$ & $0.834$ & $+0.0286$ & $-0.0002$ & $0.0006$ \\
$3.0$ & $0.743$ & $+0.0322$ & $+0.0001$ & $0.0006$ \\
\bottomrule
\end{tabular}
\caption{Bias under informative coding, $\pi = 0.12$, $200$ replicates
per row. The coding slope $b_1$ controls the C2 violation; $b_1 = 0$ is
non-informative coding. $\rho$ is the severity-coding correlation among
true cases, and is \emph{non-monotone} in $b_1$ (it rises to $0.994$ at
$b_1 = 0.5$ then falls to $0.743$ at $b_1 = 3.0$), so violation
severity is indexed by $|b_1|$, not $\rho$. The calibrated column uses
the \emph{deployable} weighted-by-$\widehat\Prob(\Y\mid S)$ estimator
(\cref{sec:methodology}; no latent cohort labels), with the same
representative review used throughout. MCSE is the Monte Carlo
standard error (sample standard deviation across replicates divided by
$\sqrt{200}$). The code-only bias is driven by the coding mechanism and
is many MCSEs from zero at every $b_1$; the deployable calibrated
estimator's bias sits within roughly two MCSEs of zero throughout,
statistically indistinguishable from unbiased. The oracle reference
(averaging over $\{S : \Y = 1\}$ in the full cohort) gives calibrated
bias within $0.003$ of the deployable column at every $b_1$, so the
remedy does not depend on knowing the latent labels.}
\label{tab:informative}
\end{table}

The correlation $\rho$ is near $1$ for small $b_1$ and declines as
$b_1$ grows, because a steep logistic in severity saturates and becomes
nonlinear, loosening the linear association between $S$ and the code
probability. The informative-coding contribution to the bias is
nonetheless real at every $b_1 > 0$: the bound \eqref{eq:meth-bound}
controls the gap between the marginal cohort sensitivity and the
average-severity coding probability by $|b_1|\,\sigma_{S,1}/4$, and the
swing in marginal sensitivity as $b_1$ grows is what drives the swing
in code-only bias.

\paragraph{The case-mix gap.} \Cref{cor:casemix} predicts that when the
chart-reviewed subsample is not representative of the cohort, the
calibrated estimator inherits a residual bias proportional to the
case-mix sensitivity gap. We test this by drawing a severity-biased
chart-review subsample: a patient's probability of being selected for
review rises with severity, controlled by a selection slope. As the
selection slope grows from $0$ to $2$, the sensitivity gap between the
reviewed subsample and the cohort grows from $0.001$ to $0.036$
(\Cref{tab:casemix}).

A calibrated estimator that marginalizes code sensitivity over the
reviewed subsample's own severity distribution inherits this gap: its
residual bias grows in magnitude from $-0.0007$ at gap $0.001$ to
$-0.006$ at gap $0.036$. A calibrated estimator that marginalizes over
the cohort severity distribution does not: its bias stays within
$\pm 0.001$ of zero throughout (\Cref{fig:informative}B). The
sensitivity gap saturates near $0.036$ for selection slopes from $1.0$
to $2.0$ because the per-case code probability $\mathrm{logit}^{-1}(b_0
+ b_1 S)$ is itself bounded above by $1$: once severity-biased
selection concentrates the reviewed subsample at the top of the
severity distribution, the subsample-marginal sensitivity approaches
the logistic ceiling and the gap stops growing. The marginal
sensitivity over true cases itself rises from $0.500$ at $b_1 = 0$
toward (but not reaching) $0.902$ at $b_1 = 3$, the same logistic
saturation. This is the phenotyping analog of the ERCC-endogenous
capture-efficiency gap of \citet{towell2026scrnacoarsening}:
chart-reviewed patients restore identifiability only to the extent
that they are drawn from the same population as the cohort being
corrected.

\begin{table}[ht]
\centering
\begin{tabular}{rrrr}
\toprule
review selection slope & sensitivity gap & subsample-ref bias & cohort-ref bias \\
\midrule
$0.0$ & $0.001$ & $-0.0007$ & $-0.0005$ \\
$0.5$ & $0.026$ & $-0.0039$ & $+0.0003$ \\
$1.0$ & $0.034$ & $-0.0046$ & $+0.0005$ \\
$1.5$ & $0.037$ & $-0.0064$ & $-0.0008$ \\
$2.0$ & $0.036$ & $-0.0060$ & $-0.0003$ \\
\bottomrule
\end{tabular}
\caption{Residual bias of the chart-review-calibrated estimator under
a severity-biased review subsample, $\pi = 0.12$, $200$ replicates per
row. The subsample-reference estimator inherits the case-mix gap; the
cohort-reference estimator does not.}
\label{tab:casemix}
\end{table}

\begin{figure}[t]
\centering
\includegraphics[width=0.95\linewidth]{informative_coding.pdf}
\caption{(A) Prevalence bias as the coding-mechanism severity slope
$b_1$ grows (the C2 violation strengthens). The code-only estimator's
bias is governed by the coding mechanism; the chart-review-calibrated
estimator stays near zero. (B) Residual bias of the calibrated
estimator as the case-mix sensitivity gap grows. Using the cohort
severity distribution as the reference removes the residual bias;
using the reviewed subsample's own distribution inherits it.}
\label{fig:informative}
\end{figure}

\subsection{Real-data plan: MIMIC-IV}

The natural real-data anchor is MIMIC-IV \citep{johnson2023mimiciv}, a
publicly available critical-care EHR dataset containing ICD diagnosis
codes, procedures, and medication orders alongside rich clinical data
(laboratory results, vital signs, clinician notes). The clinical data
can serve as a chart-review proxy: a structured rule over labs and
physiology, or a manual review of notes, supplies an adjudicated state
$\Y$ that plays the role of the chart-reviewed singleton candidate set,
while the ICD code plays the role of the coarse observation $C$.

The planned analysis takes a condition with a recognized
code-versus-chart discrepancy (for example, acute kidney injury, where
the diagnosis code is known to undercount laboratory-defined cases) and
(i) estimates code-only prevalence from the ICD code, (ii) estimates
the chart-proxy prevalence from the laboratory definition, (iii) fits
the coding model and the severity dependence on the chart-proxy
subsample, (iv) deconvolves cohort prevalence via \eqref{eq:deconvolve},
and (v) checks code-frequency consistency (\cref{thm:code-total})
directly on the cohort. For AKI we preregister the chart-proxy as
KDIGO stage 1 or higher \citep{kdigo2012aki}: a serum creatinine
increase of at least $0.3$ mg/dL within $48$ hours, an increase to at
least $1.5\times$ baseline within $7$ days, or urine output below $0.5$
mL/kg/h sustained for at least $6$ hours. A stratified chart-proxy subsample, balanced
across encounter intensity and code status, lets the case-mix gap
$\delta$ of \cref{cor:casemix} be estimated and reported. This
real-data application is described here and left as the principal
pending item; the simulation is the core validation for the present
version.
